\input{../tex-inputs/latex-leseansicht-vorspann}

\section[Arthur Schnitzler an Gustav Schwarzkopf, 13. 8. 1915]{L04490 Arthur Schnitzler an Gustav Schwarzkopf, 13. 8. 1915}
\nopagebreak\mylabel{L04490v}
\rehead{ }\normalsize\beginnumbering\briefempfaengerindex{Schwarzkopf, Gustav@\textsc{Schwarzkopf, Gustav}!zzzSchnitzler, Arthur@\emph{von Arthur Schnitzler}!1915-08-131@{13. 8. 1915}|(be}
\toendnotes[C]{\smallbreak\pagebreak[2]}
\correspDesc{Versand  durch Arthur Schnitzler am 13. 8. 1915 in Bad Ischl
\newline{}Übermittlung  am 13. 8. 1915 in Bad Ischl
\newline{}Erhalt  durch Gustav Schwarzkopf im Zeitraum [14. 8. 1915 – 18. 8. 1915?] in Wien}\toendnotes[C]{\smallbreak}
\Standort{DLA, A:Schnitzler, HS.1985.1.1897.}
\physDesc{Bildpostkarte, 523 Zeichen
\newline{}Handschrift: Bleistift, deutsche Kurrent
\newline{}Versand: Stempel: »\nobreak{}\oindex{Bad Ischl@\textbf{Bad Ischl}|pwk}Bad Ischl, 13. VIII. 15, 8\nobreak{}«.  }\toendnotes[C]{\smallbreak}
\pstart
           \noindent{}\centering{}{\pb}\textcolor{gray}{\textbf{Bad-Ischl\oindex{Bad Ischl@\textbf{Bad Ischl}|pw}.}}\pend
           
\pstart
           \centering{}\textcolor{gray}{\textbf{Postgasse\hspace*{2.5em}Theater\oindex{Lehártheater@\textbf{Lehártheater}|pw}.}}\pend
           \pstart{}{\pb}\textsc{Herrn Gustav
                     Schwarzkopf}\pend{}\pstart{}\textsc{Wien I}\oindex{I., Innere Stadt@\textbf{I., Innere Stadt}|pw}\pend{}\pstart{}\textsc{Tiefer
                        Graben 17}\oindex{Wien@\textbf{Wien}!I., Innere Stadt@\textbf{I., Innere Stadt}!Tiefer Graben 17@\textbf{Tiefer Graben 17}|pw}\pend{}{\bigskip}\vspace{1em}
\pstart
           \raggedleft{}13. 8. 15\pend
           \vspace{0.5em}
\pstart
           \textcolor{gray}{\textbf{lieber}} Guſtav, vielen Dank für Ihre
               erfreulichen \label{K_L04490-1v}\edtext{Nachrichten über Lili\pwindex{Cappellini, Lili 13.\,9.\,1909 Wien – 26.\,7.\,1928 Venedig@\textsc{Cappellini, Lili} (13.\,9.\,1909 Wien – 26.\,7.\,1928 Venedig)|pw}}{\lemma{\textnormal{\emph{Nachrichten über Lili}}}\Cendnote{\textnormal{Vgl. XXXX Auszeichnungsfehler: Dokument L04245 nicht gefunden.}}}\label{K_L04490-1};
               hoffentlich wiederholen Sie Ihren Besuch bald! Hier regnet es, wie Sie wohl wiſſen,
               allzuviel, aber man behagt{ }ſich doch ganz gut, und die Abreiſe (vielleicht \label{K_L04490-2v}\edtext{über Salzburg\oindex{Salzburg@\textbf{Salzburg}|pw}}{\lemma{\textnormal{\emph{über Salzburg}}}\Cendnote{\textnormal{Auf der Rückreise nach
                     Wien\oindex{Wien@\textbf{Wien}|pwk} gab es vom 28. 8. 1915 bis zum 30. 8. 1915 einen Aufenthalt in Salzburg\oindex{Salzburg@\textbf{Salzburg}|pwk}.}}}\label{K_L04490-2} – (Liesl\pwindex{Steinrück, Elisabeth 19.\,11.\,1885 – 7.\,4.\,1920 Partenkirchen@\textsc{Steinrück, Elisabeth} (19.\,11.\,1885 – 7.\,4.\,1920 Partenkirchen)|pw}))
               iſt noch {\pb}nicht feſtgeſetzt. Seit geſtern iſt \textsc{Mimi Giustiniani}\pwindex{Zuckerkandl, Marianne 6.\,8.\,1882 Wien – 1964 Ascona@\textsc{Zuckerkandl, Marianne} (6.\,8.\,1882 Wien – 1964 Ascona)|pw} hier; zuweilen{ }ſetzt ſich der ſehr amüſante u kluge \textsc{Friedell\pwindex{Friedell, Egon 21.\,1.\,1878 Wien – 16.\,3.\,1938 ebd.@\textsc{Friedell, Egon} (21.\,1.\,1878 Wien – 16.\,3.\,1938 ebd.)|pw}} an unſern Tiſch. Die Einakter\pwindex{Schnitzler, Arthur 15. 5. 1862 Wien – 21. 10. 1931 ebd.@\textsc{Schnitzler, Arthur} (15. 5. 1862 Wien – 21. 10. 1931 ebd.)!Komödie der Worte. Drei Einakter@\strich\emph{Komödie der Worte. Drei Einakter}|pwv} hab ich Hugo\pwindex{Hofmannsthal, Hugo von 1.\,2.\,1874 Wien – 15.\,7.\,1929 Rodaun@\textsc{Hofmannsthal, Hugo von} (1.\,2.\,1874 Wien – 15.\,7.\,1929 Rodaun)|pw}, Richard\pwindex{Beer-Hofmann, Richard 11.\,7.\,1866 Wien – 26.\,9.\,1945 New York City@\textsc{Beer-Hofmann, Richard} (11.\,7.\,1866 Wien – 26.\,9.\,1945 New York City)|pw}Kaufma{\geminationn}\pwindex{Kaufmann, Arthur 4.\,4.\,1872 Iași – 25.\,7.\,1938 Wien@\textsc{Kaufmann, Arthur} (4.\,4.\,1872 Iași – 25.\,7.\,1938 Wien)|pw} mit viel Eifer \label{K_L04490-3v}\edtext{vorgeleſen\eventindex{Villa Viola@\textbf{Villa Viola}!Private Lesung von Komödie der Worte, 11.8.1915@Private Lesung von Komödie der Worte, 11.8.1915|pwv}}{\lemma{\textnormal{\emph{vorgelesen}}}\Cendnote{\textnormal{Private Lesung von Komödie der Worte, 11.8.1915. In: A. S.: \emph{Kulturveranstaltungen}, \url{https://schnitzler-kultur.acdh.oeaw.ac.at/pmb184126.html}.}}}\label{K_L04490-3}.\pend
           
\pstart
           Wir grüßen herzlichſt{\\[\baselineskip]}Ihr \spacefill\mbox{Arthur}\pend
           \leftskip=0em{}\selectlanguage{ngerman}\endnumbering\briefempfaengerindex{Schwarzkopf, Gustav@\textsc{Schwarzkopf, Gustav}!zzzSchnitzler, Arthur@\emph{von Arthur Schnitzler}!1915-08-131@{13. 8. 1915}|)be}\mylabel{L04490h}
\begin{anhang}
\end{anhang}\newcommand{\dateiname}{L04490}\newcommand{\titel}{Arthur Schnitzler an Gustav Schwarzkopf, 13. 8. 1915}\newcommand{\editorInnen}{Herausgegeben von Jahnke, SelmaMüller, Martin Anton}\input{../tex-inputs/latex-leseansicht-abspann}
